\chapter{Display Controller and Physical HDMI Transmitter}
\label{chap:hdmi_display}

\section{Video Subsystem Overview}
To eliminate the need for external digital-to-analog converters (such as VGA DAC chips) or third-party HDMI bridge controllers, spGPU incorporates an entirely native, hardware-synthesized \textbf{HDMI Transmitter and Display Pipeline}. 

The subsystem interfaces directly with the physical differential pins of the HDMI Type-A connector on the PYNQ-Z1 board. Figure~\ref{fig:hdmi_pipeline} traces the signal path from the VGA timing generator to the differential output buffers.

\begin{figure}[htbp]
\centering
\resizebox{0.95\textwidth}{!}{
    \begin{tikzpicture}[
    >=Stealth,
    node distance=0.9cm and 1.0cm,
    every text node part/.style={align=center},
    box/.style={draw=black!80, thick, rounded corners=3pt, fill=white, inner sep=5pt, font=\small},
    pipenode/.style={box, fill=blue!10, draw=blue!70, minimum width=2.4cm},
    bus/.style={draw=black!70, line width=1.3pt, ->},
    sig/.style={draw=blue!70!black, thick, ->},
    ctl/.style={draw=red!80!black, dashed, thick, ->}
]

% ==================== PIXEL COUNTER & TIMING ====================
\node[pipenode, minimum width=3.2cm] (counter) {
    \textbf{VGA Timing Generator}\\
    \scriptsize\texttt{counter\_pixel.vhd}\\
    \scriptsize Pixel Clock: 25.0 MHz\\
    \scriptsize Horiz: 640 visible (800 total)\\
    \scriptsize Vert: 480 visible (525 total)
};

% ==================== 2X UPSCALING ====================
\node[pipenode, fill=cyan!15, draw=cyan!80, right=1.0cm of counter] (upscale) {
    \textbf{2$\times$ Upscaler Logic}\\
    \scriptsize $H_{\text{coord}} = H_{\text{vga}}[9:1]$\\
    \scriptsize $V_{\text{coord}} = V_{\text{vga}}[9:1]$\\
    \scriptsize Zero Latency / Shift
};
\draw[bus] (counter) -- node[above, font=\tiny] {$H[9:0], V[9:0]$} (upscale);

% ==================== VRAM ACCESS ====================
\node[pipenode, fill=purple!15, draw=purple!80, right=1.0cm of upscale] (vram_read) {
    \textbf{Tile VRAM Scan}\\
    \scriptsize Front Buffer Read\\
    \scriptsize Port B Latency: 1 Cycle\\
    \scriptsize Outputs: 15-bit RGB555
};
\draw[bus] (upscale) -- node[above, font=\tiny] {$X_{\text{gpu}}, Y_{\text{gpu}}$} (vram_read);

% ==================== PIPELINE ALIGNMENT ====================
\node[pipenode, fill=amber!15, draw=amber!80, below=1.0cm of counter, xshift=2.0cm, minimum width=4.5cm] (align) {
    \textbf{Pipeline Alignment Stage (1 Cycle Delay)}\\
    \scriptsize\texttt{video\_active\_pipe <= video\_active\_s}\\
    \scriptsize\texttt{h\_sync\_pipe <= h\_sync\_s}, \texttt{v\_sync\_pipe <= v\_sync\_s}\\
    \scriptsize Prevents Display Columns Distortion
};
\draw[ctl] (counter.south) |- (align.west);

% ==================== COLOR EXPANSION ====================
\node[pipenode, fill=teal!15, draw=teal!80, below=1.0cm of vram_read, xshift=-0.5cm, minimum width=3.6cm] (color_exp) {
    \textbf{Color Expander (15b $\to$ 24b)}\\
    \scriptsize $R[7:0] = R_5 \& R_5[4:2]$\\
    \scriptsize $G[7:0] = G_5 \& G_5[4:2]$\\
    \scriptsize $B[7:0] = B_5 \& B_5[4:2]$
};
\draw[bus] (vram_read.south) -- (color_exp.north);

% ==================== TMDS ENCODER ====================
\node[pipenode, fill=indigo!15, draw=indigo!80, below=1.2cm of align, xshift=2.0cm, minimum width=4.8cm] (tmds_enc) {
    \textbf{TMDS 8b/10b Encoders (\texttt{tmds\_encoder.vhd})}\\
    \scriptsize 3$\times$ Parallel Channels (Red, Green, Blue)\\
    \scriptsize XOR/XNOR Minimization + DC Disparity Tracking\\
    \scriptsize Maps Blanking Controls to TMDS Control Words
};
\draw[bus] (color_exp.south) |- (tmds_enc.east);
\draw[ctl] (align.south) -- (align.south |- tmds_enc.north);

% ==================== SERIALIZERS ====================
\node[pipenode, fill=orange!15, draw=orange!80, below=1.0cm of tmds_enc, minimum width=4.8cm] (serdes) {
    \textbf{OSERDESE2 5:1 DDR Serializers (\texttt{serializer.vhd})}\\
    \scriptsize Master/Slave Cascaded for 10-bit Serialization\\
    \scriptsize Fast Serial Clock: 125.0 MHz ($5\times$ Pixel Clock)\\
    \scriptsize Output: 4$\times$ Serial Streams (\texttt{CLK}, \texttt{D0}, \texttt{D1}, \texttt{D2})
};
\draw[bus] (tmds_enc) -- node[right, font=\tiny] {10-bit Symbols} (serdes);

% ==================== OBUFDS & HDMI ====================
\node[box, fill=slate!20, draw=black, below=1.0cm of serdes, minimum width=3.8cm] (obuf) {
    \textbf{Differential Drivers (\texttt{OBUFDS})}\\
    \scriptsize Standard: \texttt{TMDS\_33}\\
    \scriptsize Direct Physical HDMI Port (PYNQ-Z1)
};
\draw[bus] (serdes) -- (obuf);

\end{tikzpicture}

}
\caption{Complete HDMI Display and Physical Transmission Pipeline.}
\label{fig:hdmi_pipeline}
\end{figure}

\section{VGA Timing Generation (\texttt{counter\_pixel.vhd})}
The master video timing is produced by \texttt{counter\_pixel.vhd} operating at a pixel clock frequency of $25.0\,\text{MHz}$ (standard $640 \times 480$ @ $60\,\text{Hz}$ VGA progressive timing):
\begin{itemize}
    \item \textbf{Horizontal Parameters}: $640$ active display pixels, Front Porch $= 16$, HSYNC pulse $= 96$, Back Porch $= 48$, giving a total line period of $800$ pixel clock cycles ($32.0\,\mu\text{s}$, $31.25\,\text{kHz}$ line rate).
    \item \textbf{Vertical Parameters}: $480$ active display scanlines, Front Porch $= 10$, VSYNC pulse $= 2$, Back Porch $= 33$, yielding a total frame period of $525$ lines ($16.8\,\text{ms}$, $\approx 59.52\,\text{Hz}$ refresh rate).
\end{itemize}

\section{Hardware $2\times$ Nearest-Neighbor Upscaling}
The internal rendering resolution of spGPU is $320 \times 240$, whereas the HDMI transmitter requires a $640 \times 480$ raster. 

Rather than consuming memory bandwidth or DSP resources with bilinear filtering buffers, spGPU executes a zero-latency, zero-cost **$2\times$ Nearest-Neighbor Upscaling** directly in hardware by truncating the least significant bit (LSB) of the 10-bit VGA counter coordinates:
\begin{equation}
H_{\text{gpu}}[8:0] = H_{\text{vga}}[9:1] \implies H_{\text{gpu}} = \lfloor H_{\text{vga}} / 2 \rfloor
\end{equation}
\begin{equation}
V_{\text{gpu}}[8:0] = V_{\text{vga}}[9:1] \implies V_{\text{gpu}} = \lfloor V_{\text{vga}} / 2 \rfloor
\end{equation}
Each internal pixel rendered into VRAM is thus held for two consecutive pixel clocks horizontally and two consecutive scanlines vertically, creating a sharp, pixel-perfect $2 \times 2$ pixel block on the high-definition monitor.

\section{BRAM Pipeline Latency Alignment}
Synchronous Xilinx Block RAM primitives introduce an intrinsic 1-clock-cycle read latency between address presentation and valid data output. If the video synchronization pulses (\texttt{h\_sync}, \texttt{v\_sync}) and the active video flag (\texttt{video\_active}) were transmitted concurrently with address generation, the displayed image would be horizontally shifted by one pixel and the rightmost column would be truncated.

The module \texttt{vga\_struct.vhd} resolves this by routing the timing flags through dedicated 1-cycle pipeline registers:
\begin{lstlisting}[language=VHDL, caption={Pipeline Latency Alignment in \texttt{vga\_struct.vhd}.}, label={lst:vga_pipeline}]
align_proc : process(pixelclock, rst)
begin
    if rst = '1' then
        video_active_pipe <= '0';
        h_sync_pipe       <= '1';
        v_sync_pipe       <= '1';
    elsif rising_edge(pixelclock) then
        video_active_pipe <= video_active_s;
        h_sync_pipe       <= h_sync_s;
        v_sync_pipe       <= v_sync_s;
    end if;
end process;
\end{lstlisting}
This guarantees that pixel color data leaving the BRAM arrives at the TMDS encoder in exact temporal alignment with its corresponding active video framing.

\section{Color Expansion (RGB555 to RGB888)}
The 15-bit color stored in VRAM allocates 5 bits per primary color. To drive standard 24-bit TrueColor HDMI sinks without darkening the dynamic range, the 5-bit color components are expanded to 8 bits by replicating the 3 most significant bits into the 3 least significant bit positions:
\begin{align}
R_{8}[7:0] &= R_{5}[4:0] \mathbin{\Vert} R_{5}[4:2] \\
G_{8}[7:0] &= G_{5}[4:0] \mathbin{\Vert} G_{5}[4:2] \\
B_{8}[7:0] &= B_{5}[4:0] \mathbin{\Vert} B_{5}[4:2]
\end{align}
This ensures that maximum 5-bit value (\texttt{"11111"}) maps accurately to maximum 8-bit value (\texttt{"11111111"} = 255), maintaining full dynamic contrast and color saturation.

\section{TMDS 8b/10b Encoders (\texttt{tmds\_encoder.vhd})}
Transition-Minimized Differential Signaling (TMDS) encodes 8-bit color data into 10-bit symbols to minimize high-frequency electromagnetic interference (EMI) and maintain DC balance over the physical cable:
\begin{enumerate}
    \item \textbf{Transition Minimization}: Evaluates the number of '1' bits in the input byte. Depending on the density, it generates an 8-bit sequence using either cumulative XOR or cumulative XNOR operations, asserting the 9th bit (\texttt{q\_m[8]}) to denote the chosen scheme.
    \item \textbf{DC Balance \& Disparity Tracking}: An internal signed 4-bit accumulator (\texttt{disparity}) monitors the running discrepancy between transmitted '1's and '0's. If the disparity exceeds zero, the 10th bit is inverted to restore the DC baseline of the transmission line.
    \item \textbf{Blanking Control Word Mapping}: During the horizontal and vertical blanking intervals (\texttt{video\_active = '0'}), the encoder maps the 2-bit control inputs (\texttt{h\_sync}, \texttt{v\_sync}) to one of four unique 10-bit synchronization sequences (\texttt{"1101010100"}, \texttt{"0010101011"}, etc.) exhibiting high transition rates for phase-locked loop (PLL) clock recovery at the sink.
\end{enumerate}

\section{Physical Serialization and Differential I/O (\texttt{serializer.vhd})}
To transmit 10-bit parallel TMDS symbols over a single-conductor pair per channel at $60\,\text{Hz}$, the bit rate per lane is:
\begin{equation}
\text{Bit Rate} = 25.0\,\text{MHz} \times 10 = 250.0\,\text{Mbps/lane}
\end{equation}

In Xilinx 7-series FPGAs, this is realized by cascading two \textbf{\texttt{OSERDESE2}} hardware primitives in Master/Slave configuration:
\begin{itemize}
    \item \textbf{Clocking}: Operating in Double Data Rate (DDR) mode with a 5:1 serialization ratio, the blocks are clocked by the fast serial clock \texttt{serialclock = 125.0 MHz} ($5 \times 25.0\,\text{MHz}$) and divided parallel clock \texttt{pixelclock = 25.0 MHz}.
    \item \textbf{Cascading}: The Master OSERDESE2 serializes the lower 8 bits, while the Slave OSERDESE2 supplies the upper 2 bits via internal shift lines (\texttt{SHIFTIN1}, \texttt{SHIFTIN2}).
    \item \textbf{Differential Drivers}: The serialized output bitstreams for the clock channel and the three data channels (Red, Green, Blue) drive four Xilinx \textbf{\texttt{OBUFDS}} differential output buffer primitives configured with the \texttt{TMDS\_33} I/O standard, providing the correct current and impedance matching for direct connection to any standard HDMI monitor.
\end{itemize}
